\section{Method}
\label{sec:method}

\Cref{fig:architecture} gives the full picture.
Frames arrive one at a time from a live video stream. At each step $t$, the Semantic Keyframe Memory (\cref{sec:skm}) scores the new frame $I_t$ and decides whether to store it in a bank of at most $N$ entries.
Independently, a user may submit a question $q$ at any moment.
The Temporal Query Router (\cref{sec:tqr}) classifies $q$ into a temporal scope and passes only the matching memory entries forward.
The Stream-Fused Generator (\cref{sec:sfg}) then runs a two-stage perception-synthesis pipeline over those entries and returns an answer in under 200\,ms.

\subsection{Semantic Keyframe Memory}
\label{sec:skm}

A live video stream produces frames at 15 to 30 fps. Consecutive frames are redundant.
Storing all frames is not feasible. Uniform subsampling loses rare but informative events.
The SKM retains frames based on visual-semantic importance.

\paragraph{Visual encoding.}
Each incoming frame $I_t$ passes through a CLIP ViT-B/32~\cite{radford2021clip} encoder, producing a 512-dimensional embedding $v_t$.
Not every frame goes through encoding. By default the system processes every 5th frame, giving 3 to 6 candidates per second and keeping encoding cost low.
The system saves a small JPEG thumbnail alongside each embedding for later visualization.

\paragraph{Importance scoring.}
The system scores every candidate against what is already in memory. Two signals determine the score.

\textbf{Visual novelty} captures how different the candidate looks. We compute the cosine similarity between $v_t$ and each stored embedding, and one minus the maximum gives the novelty. If every stored frame looks similar to the candidate, novelty is low. If no stored frame is close, novelty is high.

\textbf{Temporal coverage} captures how well the candidate fills a gap in the timeline. It is the smallest absolute time difference between $t$ and any stored timestamp, divided by a horizon $T_{\max}$.

The importance score for frame $I_t$ is:
\begin{equation}
\begin{split}
  s(I_t) =\; & \alpha \Bigl(1 - \max_{j \in \mathcal{M}} \cos(v_t, v_j)\Bigr) \\
  & + (1{-}\alpha)\,\min\!\Bigl(\frac{\min_{j \in \mathcal{M}} |t - t_j|}{T_{\max}},\, 1\Bigr)
\end{split}
  \label{eq:importance}
\end{equation}
where $\mathcal{M}$ is the current memory set, $\alpha = 0.7$ controls the novelty-coverage trade-off, and $T_{\max} = 300$\,s (5 minutes). A frame that is both visually distinct and temporally distant from stored entries receives the highest score.

\paragraph{Memory update.}
Below capacity, every candidate enters memory directly.
Once memory is full, the system re-scores stored entries and finds the lowest-scoring one. If the incoming candidate scores higher, it replaces that entry. If not, the system drops the candidate. The full update runs in $O(N)$ time and adds under 4\,ms of latency.

\paragraph{Temporal bookkeeping.}
Each memory entry stores its original timestamp. The TQR uses these timestamps to filter entries by time window. We sort memory by timestamp after each insertion.


\subsection{Temporal Query Router}
\label{sec:tqr}

When a question $q$ arrives, the TQR determines which part of memory is relevant.
\Cref{fig:tqr} shows the pipeline. The system matches the question against keyword lists, the scope selector picks the best temporal scope, and the memory filter returns entries in that time window.

\paragraph{Scope classification.}
The TQR assigns each question one of three temporal scopes: \textsc{instant}, \textsc{recent}, or \textsc{historical}.
A keyword matcher runs over the lowercased query string.
Three manually curated keyword lists drive the decision. \textsc{Instant} cues include ``right now,'' ``currently,'' and ``at this moment.'' \textsc{Recent} cues include ``just,'' ``a moment ago,'' ``few seconds,'' and ``last minute.'' \textsc{Historical} cues include ``earlier,'' ``before,'' ``previously,'' and ``how many times.''

The scope with the most keyword hits wins. When no keyword matches or two scopes tie with low confidence (top score / total $<$ 0.6), the system defaults to \textsc{historical}, granting full memory access. This fallback is conservative: an over-broad context is better than a truncated one.

\paragraph{Scope-aware filtering.}
Given predicted scope $c$ and current time $t_{\text{now}}$, the filtered subset $\mathcal{F} \subseteq \mathcal{M}$ is:
\begin{equation}
  \mathcal{F} \!=\!
  \begin{cases}
    \{\argmax_{j \in \mathcal{M}} t_j\}                        & c = \text{inst.} \\[3pt]
    \{j \!\in\! \mathcal{M} : t_j \!\ge\! t_{\text{now}} - W\} & c = \text{rec.}  \\[3pt]
    \mathcal{M}                                                 & c = \text{hist.}
  \end{cases}
  \label{eq:scope_filter}
\end{equation}
where $W = 30$\,s is the recency window. \textsc{inst.}, \textsc{rec.}, and \textsc{hist.}\ stand for the three scopes defined above. Only entries in $\mathcal{F}$ reach the Stream-Fused Generator.


\subsection{Stream-Fused Generator}
\label{sec:sfg}

The SFG turns filtered keyframes into a natural-language answer through two stages.

\paragraph{Stage 1: BLIP visual perception.}
The system samples up to 6 frames from the filtered set (or takes all of them if fewer than 6 pass the scope filter).
Two BLIP~\cite{li2022blip} models process each frame independently:

\begin{itemize}[itemsep=2pt]
  \item \textbf{BLIP captioning} (blip-image-captioning-base): generates a free-form description of the frame contents. Each caption is at most 50 tokens.
  \item \textbf{BLIP VQA} (blip-vqa-base): takes the user's question and the frame as input. It produces a short direct answer for that specific frame. Each answer is at most 30 tokens.
\end{itemize}

The system collects all captions and VQA answers into a pool of textual observations.

\paragraph{Stage 2: Flan-T5 language synthesis.}
The system assembles a structured prompt from the pool. It contains the user's question, unique scene descriptions (deduplicated Stage-1 captions), direct frame-level answers (deduplicated Stage-1 VQA outputs), and the count of sampled frames.

Flan-T5-base~\cite{chung2022flanT5} receives this prompt and generates up to 80 tokens with beam search (4 beams, length penalty 1.0, no-repeat 3-gram constraint). The result is one coherent answer that ties together observations from multiple keyframes.

When Flan-T5 is unavailable (e.g., CPU-only deployments), the system falls back to template concatenation of deduplicated VQA answers and captions.

\paragraph{Latency budget.}
The total per-query latency breaks down as follows. TQR scope classification takes about 1\,ms. Frame sampling and BLIP inference on 6 frames takes 80 to 120\,ms. Flan-T5 decoding takes 60 to 80\,ms.
End-to-end latency stays under 200\,ms on a single NVIDIA A100.


\subsection{Training}
\label{sec:training}

Every model in the pipeline is a frozen pre-trained checkpoint. The system requires no end-to-end training.

\paragraph{CLIP encoder.}
The vision encoder is the CLIP ViT-B/32~\cite{radford2021clip} checkpoint from OpenAI. It takes 224$\times$224 images and outputs 512-dimensional embeddings. We freeze all weights.

\paragraph{BLIP perception models.}
Two BLIP~\cite{li2022blip} checkpoints from Salesforce handle perception: blip-image-captioning-base for captioning and blip-vqa-base for visual question answering. We load both in evaluation mode with no gradient computation.

\paragraph{Flan-T5 synthesis.}
The language synthesis model is Flan-T5-base~\cite{chung2022flanT5} from Google, instruction-tuned on 1,836 tasks. It reads the structured Stage-2 prompts and generates answers. We apply no fine-tuning.

\paragraph{SKM hyperparameters.}
Because the importance function operates on frozen CLIP features, it needs no learned parameters.
We selected the novelty weight $\alpha$ and temporal horizon $T_{\max}$ by sweeping $\alpha \in \{0.5, 0.6, 0.7, 0.8\}$ and $T_{\max} \in \{60, 120, 300, 600\}$ on a held-out validation set. We chose the best combination ($\alpha = 0.7$, $T_{\max} = 300$\,s) by LiveQA-Bench accuracy.

\paragraph{TQR keyword lists.}
We compiled the keyword lists by hand from annotated question-scope pairs in the LiveQA-Bench training split. We collected common temporal expressions for each scope (e.g., ``right now'' for instant, ``earlier'' for historical) and validated them on a held-out split. The matcher classifies questions with clear temporal cues correctly in most cases. Ambiguous questions default to the historical scope, giving the generator full memory access.
